\documentclass[10pt,a4paper]{amsart}
\usepackage{amsmath,amssymb,amsthm,graphicx}
% For commutative diagrams we use Xy-pic
\usepackage[all]{xy}

% We use 2cell for 2-commutative diagrams.
\xyoption{2cell}
% Theorem environments.
%
\theoremstyle{plain}
\newtheorem{theorem}[subsection]{Theorem}
\newtheorem{proposition}[subsection]{Proposition}
\newtheorem{lemma}[subsection]{Lemma}

\theoremstyle{definition}
\newtheorem{definition}[subsection]{Definition}
\newtheorem{example}[subsection]{Example}
\newtheorem{exercise}[subsection]{Exercise}
\newtheorem{situation}[subsection]{Situation}

\theoremstyle{remark}
\newtheorem{remark}[subsection]{Remark}
\newtheorem{remarks}[subsection]{Remarks}

\numberwithin{equation}{subsection}

% Macros
%
\def\lim{\mathop{\mathrm{lim}}\nolimits}
\def\colim{\mathop{\mathrm{colim}}\nolimits}
\def\Spec{\mathop{\mathrm{Spec}}}
\def\Hom{\mathop{\mathrm{Hom}}\nolimits}
\def\Ext{\mathop{\mathrm{Ext}}\nolimits}
\def\SheafHom{\mathop{\mathcal{H}\!\mathit{om}}\nolimits}
\def\SheafExt{\mathop{\mathcal{E}\!\mathit{xt}}\nolimits}
\def\Sch{\mathit{Sch}}
\def\Mor{\mathop{\mathrm{Mor}}\nolimits}
\def\Ob{\mathop{\mathrm{Ob}}\nolimits}
\def\Sh{\mathop{\mathit{Sh}}\nolimits}
\def\NL{\mathop{N\!L}\nolimits}
\def\CH{\mathop{\mathrm{CH}}\nolimits}
\def\proetale{{pro\text{-}\acute{e}tale}}
\def\etale{{\acute{e}tale}}
\def\QCoh{\mathit{QCoh}}
\def\Ker{\mathop{\mathrm{Ker}}}
\def\Im{\mathop{\mathrm{Im}}}
\def\Coker{\mathop{\mathrm{Coker}}}
\def\Coim{\mathop{\mathrm{Coim}}}

% Boxtimes
%
\DeclareMathSymbol{\boxtimes}{\mathbin}{AMSa}{"02}

%
% Macros for moduli stacks/spaces
%
\def\QCohstack{\mathcal{QC}\!\mathit{oh}}
\def\Cohstack{\mathcal{C}\!\mathit{oh}}
\def\Spacesstack{\mathcal{S}\!\mathit{paces}}
\def\Quotfunctor{\mathrm{Quot}}
\def\Hilbfunctor{\mathrm{Hilb}}
\def\Curvesstack{\mathcal{C}\!\mathit{urves}}
\def\Polarizedstack{\mathcal{P}\!\mathit{olarized}}
\def\Complexesstack{\mathcal{C}\!\mathit{omplexes}}
% \Pic is the operator that assigns to X its picard group, usage \Pic(X)
% \Picardstack_{X/B} denotes the Picard stack of X over B
% \Picardfunctor_{X/B} denotes the Picard functor of X over B
\def\Pic{\mathop{\mathrm{Pic}}\nolimits}
\def\Picardstack{\mathcal{P}\!\mathit{ic}}
\def\Picardfunctor{\mathrm{Pic}}
\def\Deformationcategory{\mathcal{D}\!\mathit{ef}}
\usepackage{cleveref}
\crefname{originallemma}{original lemma}{original lemmas}
\Crefname{originallemma}{Original lemma}{Original lemmas}
\makeatletter
\newlabel{coherent-lemma-induction-principle}{{08DR}{source}{}{sourcepointer-1}{}}
\newlabel{coherent-lemma-induction-principle@cref}{{[originallemma][0][]08DR}{[1][0][]source}{}{}{}}
\makeatother
\crefname{originalreference}{original reference}{original references}
\Crefname{originalreference}{Original reference}{Original references}
\makeatletter
\newlabel{lemma-interlude-I}{{0F0X}{source}{}{sourcepointer-2}{}}
\newlabel{lemma-interlude-I@cref}{{[originalreference][0][]0F0X}{[1][0][]source}{}{}{}}
\makeatother
\crefname{originalreference}{original reference}{original references}
\Crefname{originalreference}{Original reference}{Original references}
\makeatletter
\newlabel{lemma-mayer-vietoris}{{01EB}{source}{}{sourcepointer-3}{}}
\newlabel{lemma-mayer-vietoris@cref}{{[originalreference][0][]01EB}{[1][0][]source}{}{}{}}
\makeatother
\crefname{originalreference}{original reference}{original references}
\Crefname{originalreference}{Original reference}{Original references}
\makeatletter
\newlabel{lemma-relative-mayer-vietoris}{{01EC}{source}{}{sourcepointer-4}{}}
\newlabel{lemma-relative-mayer-vietoris@cref}{{[originalreference][0][]01EC}{[1][0][]source}{}{}{}}
\makeatother
\crefname{originalreference}{original reference}{original references}
\Crefname{originalreference}{Original reference}{Original references}
\makeatletter
\newlabel{proposition-cd-affine}{{0F0V}{source}{}{sourcepointer-5}{}}
\newlabel{proposition-cd-affine@cref}{{[originalreference][0][]0F0V}{[1][0][]source}{}{}{}}
\makeatother
\title{Finiteness of the cohomological dimension of finite type morphisms and finite type schemes over a field}
\author{The Stacks Project Authors}
\date{stacks-project:etale-cohomology.tex:local-snapshot-20261011; GNU FDL 1.2 or later}
\begin{document}
\maketitle
\noindent\textit{Editorial note.} All delivered mathematical text is the author's own, copied from the pinned Stacks Project chapter file etale-cohomology.tex; no statement or proof step was rewritten, reordered or removed.\par\medskip
\noindent\textit{Editorial note.} The counted proof is the author's complete argument: the two applications of the induction principle, the affine base cases and the Mayer-Vietoris steps propagating the bound from the pieces to their union.\par\medskip
\noindent\textit{Editorial note.} The label lemma-finite-cd names the counted statement, and the proof environment immediately following it in the source is delivered verbatim as its proof.\par\medskip
\noindent\textit{Editorial note.} The counted statement is carried by the author's own lemma environment in the chapter source.\par\medskip
\noindent\textit{Editorial note.} Pre-open gate: python operations/agent-authoring/tools/precard\_gate.py --source etale-cohomology.tex --label lemma-finite-cd --start 20040 --end 20071 --local-id 3566 returned RESULT=OK (L3566 free, occupancy FREE, public interval clear). Candidate taken from lane-s4/halfB-priority.json.\par\medskip
\noindent\textit{Editorial note.} The statement and proof are copied from the pinned chapter file; the cited inputs are kept as external pointers with their canonical Stacks tags.\par\medskip
\noindent\textit{Allowed standard prerequisites (editorial).}\begin{itemize}
\item The induction principle for Noetherian schemes. Cohomology of Schemes, the lemma allowing a statement about Noetherian schemes to be proved by checking affine schemes and the Mayer-Vietoris step for X = U union V. Supplies the induction pattern used twice, reducing the assertions to the affine case and to the case X = U union V.
\item Finiteness of the cohomological dimension in the affine case. Same chapter, the lemma bounding the cohomological dimension of a finite type morphism with affine target, and the proposition bounding the cohomological dimension of an affine scheme over a field of finite cohomological dimension. Supply the base cases of the two inductions.
\item Relative and absolute Mayer-Vietoris sequences. Same chapter, the relative Mayer-Vietoris sequence for a morphism and the Mayer-Vietoris sequence for the cohomology of a scheme. Supply the induction steps, since a bound valid for U, V and U cap V propagates to X = U union V.
\end{itemize}\medskip
\ifdefined\hypertarget\hypertarget{sourcepointer-1}{}\fi
\noindent\textit{Source reference.} Original lemma 08DR: sources/stacks-chow/tags.txt, source lines 10646--10646. Prerequisite pointer invoked at proof line(s) 20054; it names a standard fact of a sibling chapter that the counted argument uses, and every step of the counted proof built on it is retained verbatim in the delivered body.\par\medskip
\ifdefined\hypertarget\hypertarget{sourcepointer-2}{}\fi
\noindent\textit{Source reference.} Original reference 0F0X: sources/stacks-etale-cohomology/pinned-original-source.tex, source lines 19941--19941. Prerequisite pointer invoked at proof line(s) 20056; it names a standard fact of the same chapter that the counted argument uses, and every step of the counted proof built on it is retained verbatim in the delivered body.\par\medskip
\ifdefined\hypertarget\hypertarget{sourcepointer-3}{}\fi
\noindent\textit{Source reference.} Original reference 01EB: sources/stacks-etale-cohomology/pinned-original-source.tex, source lines 7176--7176. Prerequisite pointer invoked at proof line(s) 20070; it names a standard fact of the same chapter that the counted argument uses, and every step of the counted proof built on it is retained verbatim in the delivered body.\par\medskip
\ifdefined\hypertarget\hypertarget{sourcepointer-4}{}\fi
\noindent\textit{Source reference.} Original reference 01EC: sources/stacks-etale-cohomology/pinned-original-source.tex, source lines 7223--7223. Prerequisite pointer invoked at proof line(s) 20060; it names a standard fact of the same chapter that the counted argument uses, and every step of the counted proof built on it is retained verbatim in the delivered body.\par\medskip
\ifdefined\hypertarget\hypertarget{sourcepointer-5}{}\fi
\noindent\textit{Source reference.} Original reference 0F0V: sources/stacks-etale-cohomology/pinned-original-source.tex, source lines 19763--19763. Prerequisite pointer invoked at proof line(s) 20066; it names a standard fact of the same chapter that the counted argument uses, and every step of the counted proof built on it is retained verbatim in the delivered body.\par\medskip
\begin{lemma}
\label{lemma-finite-cd}
Let $K$ be a field.
\begin{enumerate}
\item If $f : X \to Y$ is a morphism of finite type schemes over $K$,
then $\text{cd}(f) < \infty$.
\item If $\text{cd}(K) < \infty$, then $\text{cd}(X) < \infty$
for any finite type scheme $X$ over $K$.
\end{enumerate}
\end{lemma}

\begin{proof}
Proof of (1). We may assume $Y$ is affine. We will use the
induction principle of 
Cohomology of Schemes, Lemma \ref{coherent-lemma-induction-principle}
to prove this. If $X$ is affine too, then the result holds by
Lemma \ref{lemma-interlude-I}. Thus it suffices to show that if
$X = U \cup V$ and the result is true for $U \to Y$, $V \to Y$, and
$U \cap V \to Y$, then it is true for $f$. This follows from the
relative Mayer-Vietoris sequence, see
Lemma \ref{lemma-relative-mayer-vietoris}.

\medskip\noindent
Proof of (2). We will use the induction principle of
Cohomology of Schemes, Lemma \ref{coherent-lemma-induction-principle}
to prove this. If $X$ is affine, then the result holds by
Proposition \ref{proposition-cd-affine}. Thus it suffices to show that if
$X = U \cup V$ and the result is true for $U$, $V$, and
$U \cap V $, then it is true for $X$. This follows from the
Mayer-Vietoris sequence, see
Lemma \ref{lemma-mayer-vietoris}.
\end{proof}

\end{document}
